%! Introducing Vectors — Unit 1, Lesson 1.0 — Guided Notes (Answer Key)
\documentclass[10pt]{article}
\usepackage{linalg-article}
\usepackage{linalg-key}   % pulls in -boxes; do NOT also load -boxes
\newcommand{\TallMath}[1]{$\displaystyle #1\rule[-1.4em]{0pt}{3.2em}$}
\newcommand{\vv}{\mathbf{v}}
\newcommand{\ww}{\mathbf{w}}

\begin{document}

\pageheader{Unit 1, Lesson 1.0}{Guided Notes --- Answer Key}

\vspace{0.12in}

\begin{objectivebox}
\textbf{By the end of this lesson, I will be able to\dots}
\begin{itemize}\setlength{\itemsep}{2pt}
  \item describe a vector by its \emph{components} and as an \emph{arrow} in the plane;
  \item add two vectors and multiply a vector by a scalar;
  \item find the \emph{length} (magnitude) of a vector and explain what it measures.
\end{itemize}
\end{objectivebox}

\vspace{0.10in}

\begin{vocabbox}
\small
Fill in each term as we build it together.
\par\vspace{2pt}
\vocabans{Vector}{an ordered list of numbers; in 2-D, two components picturing an arrow (or point) in the plane}
\vocabans{Scalar}{a single number --- it has size but no direction, and it scales a vector}
\vocabans{Components}{the individual numbers $v_1, v_2$ inside a vector; its coordinates}
\vocabans{Vector addition}{add matching components; geometrically, place the arrows tip to tail}
\vocabans{Scalar multiplication}{multiply every component by the same number $c$; this stretches, shrinks, or flips the arrow}
\vocabans{Magnitude (length)}{the length of the arrow, $\|\mathbf{v}\| = \sqrt{v_1^2 + v_2^2}$; never negative}
\vocabans{Unit vector}{a vector whose length equals $1$}
\end{vocabbox}

\begin{hookbox}
A drone lifts off and flies \textbf{3 blocks east and 4 blocks north}. A single pair
of numbers, $(3,4)$, records the \emph{whole trip}: how far east, and how far north.
\begin{itemize}\setlength{\itemsep}{1pt}
  \item What one number would tell you the drone's \emph{straight-line distance} from
        where it started? \ans{the length of the arrow, $5$ blocks}
  \item How could you find it from the $3$ and the $4$? \ansline{Pythagorean theorem: $\sqrt{3^2+4^2}=\sqrt{25}=5$}
\end{itemize}
\end{hookbox}

\begin{notesbox}{1. What Is a Vector?}
A \textbf{vector} is an ordered list of numbers. In two dimensions it has two
\textbf{components} and we write it as a column:
\[
  \vv = \begin{bmatrix} v_1 \\ v_2 \end{bmatrix},
  \qquad\text{for example}\qquad
  \vv = \begin{bmatrix} 3 \\ 4 \end{bmatrix}.
\]

We can picture $\vv$ two ways:
\begin{itemize}\setlength{\itemsep}{2pt}
  \item as a \textbf{point} in the plane at $(v_1, v_2)$, or
  \item as an \textbf{arrow} from the origin out to that point.
\end{itemize}

The first component says how far to move \ans{east (right)}; the second says how far to move
\ans{north (up)}. A single number, like $5$, is called a \textbf{scalar} --- it has size but
no direction.

\vspace{4pt}
\begin{center}
\begin{tikzpicture}[scale=0.62]
  \draw[step=1, linegray!60, thin] (-0.4,-0.4) grid (5.4,5.4);
  \draw[->, thick, charcoal] (0,0)--(5.6,0) node[right]{\scriptsize east};
  \draw[->, thick, charcoal] (0,0)--(0,5.6) node[above]{\scriptsize north};
  \foreach \x in {1,2,3,4,5} \node[below, font=\scriptsize] at (\x,-0.1){\x};
  \foreach \y in {1,2,3,4,5} \node[left, font=\scriptsize] at (-0.1,\y){\y};
  \draw[->, very thick, burgundy] (0,0)--(3,4)
        node[midway, above left]{$\vv$};
  \fill[burgundy] (3,4) circle (2.4pt) node[above right]{\small $(3,4)$};
  \draw[dashed, slate] (3,0)--(3,4);
  \draw[dashed, slate] (0,4)--(3,4);
  \node[below, slate, font=\scriptsize] at (1.5,0){$3$};
  \node[left, slate, font=\scriptsize] at (0,2){$4$};
\end{tikzpicture}
\end{center}
\end{notesbox}

\begin{notesbox}{2. Adding and Scaling Vectors}
\textbf{Vector addition} combines two vectors component by component:
\[
  \vv + \ww =
  \begin{bmatrix} v_1 \\ v_2 \end{bmatrix} +
  \begin{bmatrix} w_1 \\ w_2 \end{bmatrix} =
  \begin{bmatrix} v_1 + w_1 \\ v_2 + w_2 \end{bmatrix}.
\]
Geometrically you place the arrows \textbf{tip to tail}: walk along $\vv$, then along
$\ww$, and $\vv+\ww$ is the arrow from your start to your finish.

\vspace{2pt}
\textit{Example.} With $\vv = \begin{bmatrix} 3 \\ 4 \end{bmatrix}$ and
$\ww = \begin{bmatrix} 1 \\ -2 \end{bmatrix}$:
\qquad $\vv + \ww = \begin{bmatrix}{\color{keyred}\mathbf{4}}\\[2pt]{\color{keyred}\mathbf{2}}\end{bmatrix}$

\vspace{6pt}
\textbf{Scalar multiplication} scales every component by the same number $c$:
\[
  c\,\vv = \begin{bmatrix} c\,v_1 \\ c\,v_2 \end{bmatrix}.
\]
\begin{itemize}\setlength{\itemsep}{1pt}
  \item $c>1$ makes the arrow \ans{longer}; \; $0<c<1$ makes it \ans{shorter}.
  \item A \emph{negative} $c$ \ans{reverses} the arrow's direction.
\end{itemize}
\textit{Example.} $2\vv = \begin{bmatrix}{\color{keyred}\mathbf{6}}\\[2pt]{\color{keyred}\mathbf{8}}\end{bmatrix}$
\qquad
$-\vv = \begin{bmatrix}{\color{keyred}\mathbf{-3}}\\[2pt]{\color{keyred}\mathbf{-4}}\end{bmatrix}$

\vspace{4pt}
An expression like $c\,\vv + d\,\ww$ --- scaling and then adding --- is called a
\textbf{linear combination}. It is the central idea of this whole unit.
\end{notesbox}

\boxguard
\begin{notesbox}{3. The Length of a Vector}
The \textbf{magnitude} (or \textbf{length}) of a vector is how long its arrow is. Because
the components form the legs of a right triangle, the Pythagorean theorem gives
\[
  \|\vv\| = \sqrt{\,v_1^{\,2} + v_2^{\,2}\,}.
\]
The double bars $\|\vv\|$ are read ``the length of $\vv$.'' The length is always a
\textbf{scalar}, and it is never negative.

\vspace{2pt}
\textit{Example.} For $\vv = \begin{bmatrix} 3 \\ 4 \end{bmatrix}$:
\[
  \|\vv\| = \sqrt{3^2 + 4^2} = \sqrt{{\color{keyred}\mathbf{25}}} = {\color{keyred}\mathbf{5}}.
\]
So the drone from the Hook ends up \ans{$5$} blocks in a straight line from the start.

\vspace{4pt}
A vector whose length is exactly $1$ is called a \textbf{unit vector}. To shrink any
vector to length $1$, divide it by its own length: $\dfrac{1}{\|\vv\|}\,\vv$. This keeps
the \ans{direction} the same but sets the length to \ans{$1$}.
\end{notesbox}

\boxguard
\begin{practicebox}
Let $\vv = \begin{bmatrix} 6 \\ 8 \end{bmatrix}$ and $\ww = \begin{bmatrix} -3 \\ 0 \end{bmatrix}$.
\begin{enumerate}[leftmargin=1.6em]
  \item Find $\vv + \ww$. \hfill \ans{$\begin{bmatrix} 3 \\ 8 \end{bmatrix}$}
  \item Find $\tfrac{1}{2}\vv$. \hfill \ans{$\begin{bmatrix} 3 \\ 4 \end{bmatrix}$}
  \item Find $\|\vv\|$. Show the Pythagorean step. \ansline{$\sqrt{6^2+8^2}=\sqrt{100}=10$}
  \item Is $\ww$ a unit vector? How do you know? \ansline{No --- $\|\ww\|=\sqrt{(-3)^2+0^2}=3\neq 1$}
\end{enumerate}
\end{practicebox}

\end{document}
